\documentclass{amsart}
% For dealing with references we use the comment environment
\usepackage{verbatim}
\newenvironment{reference}{\comment}{\endcomment}
%\newenvironment{reference}{}{}
\newenvironment{slogan}{\comment}{\endcomment}
\newenvironment{history}{\comment}{\endcomment}

% For commutative diagrams we use Xy-pic
\usepackage[all]{xy}

% We use 2cell for 2-commutative diagrams.
\xyoption{2cell}
\UseAllTwocells

% We use multicol for the list of chapters between chapters
\usepackage{multicol}

% This is generally recommended for better output
\usepackage{lmodern}
\usepackage[T1]{fontenc}

% For cross-file-references

% Package for hypertext links:
\usepackage{hyperref}

% For any local file, say "hello.tex" you want to link to please

% Theorem environments.
%
\theoremstyle{plain}
\newtheorem{theorem}[subsection]{Theorem}
\newtheorem{proposition}[subsection]{Proposition}
\newtheorem{lemma}[subsection]{Lemma}

\theoremstyle{definition}
\newtheorem{definition}[subsection]{Definition}
\newtheorem{example}[subsection]{Example}
\newtheorem{exercise}[subsection]{Exercise}
\newtheorem{situation}[subsection]{Situation}

\theoremstyle{remark}
\newtheorem{remark}[subsection]{Remark}
\newtheorem{remarks}[subsection]{Remarks}

\numberwithin{equation}{subsection}

% Macros
%
\def\lim{\mathop{\mathrm{lim}}\nolimits}
\def\colim{\mathop{\mathrm{colim}}\nolimits}
\def\Spec{\mathop{\mathrm{Spec}}}
\def\Hom{\mathop{\mathrm{Hom}}\nolimits}
\def\Ext{\mathop{\mathrm{Ext}}\nolimits}
\def\SheafHom{\mathop{\mathcal{H}\!\mathit{om}}\nolimits}
\def\SheafExt{\mathop{\mathcal{E}\!\mathit{xt}}\nolimits}
\def\Sch{\mathit{Sch}}
\def\Mor{\mathop{\mathrm{Mor}}\nolimits}
\def\Ob{\mathop{\mathrm{Ob}}\nolimits}
\def\Sh{\mathop{\mathit{Sh}}\nolimits}
\def\NL{\mathop{N\!L}\nolimits}
\def\CH{\mathop{\mathrm{CH}}\nolimits}
\def\proetale{{pro\text{-}\acute{e}tale}}
\def\etale{{\acute{e}tale}}
\def\QCoh{\mathit{QCoh}}
\def\Ker{\mathop{\mathrm{Ker}}}
\def\Im{\mathop{\mathrm{Im}}}
\def\Coker{\mathop{\mathrm{Coker}}}
\def\Coim{\mathop{\mathrm{Coim}}}

% Boxtimes
%
\DeclareMathSymbol{\boxtimes}{\mathbin}{AMSa}{"02}

%
% Macros for moduli stacks/spaces
%
\def\QCohstack{\mathcal{QC}\!\mathit{oh}}
\def\Cohstack{\mathcal{C}\!\mathit{oh}}
\def\Spacesstack{\mathcal{S}\!\mathit{paces}}
\def\Quotfunctor{\mathrm{Quot}}
\def\Hilbfunctor{\mathrm{Hilb}}
\def\Curvesstack{\mathcal{C}\!\mathit{urves}}
\def\Polarizedstack{\mathcal{P}\!\mathit{olarized}}
\def\Complexesstack{\mathcal{C}\!\mathit{omplexes}}
% \Pic is the operator that assigns to X its picard group, usage \Pic(X)
% \Picardstack_{X/B} denotes the Picard stack of X over B
% \Picardfunctor_{X/B} denotes the Picard functor of X over B
\def\Pic{\mathop{\mathrm{Pic}}\nolimits}
\def\Picardstack{\mathcal{P}\!\mathit{ic}}
\def\Picardfunctor{\mathrm{Pic}}
\def\Deformationcategory{\mathcal{D}\!\mathit{ef}}

\setlength{\emergencystretch}{3em}
\begin{document}
\title[Stacks Project, Tag 062I]{340 --- First Koszul homology regularity implies quasi-regularity}
\maketitle
\noindent Copyright (C) 2005--2025 Johan de Jong. Stacks Project authors. Source: \href{https://stacks.math.columbia.edu/tag/062I}{Tag 062I}.

\noindent Editorial difficulty note (not author text). Difficulty H2: The proof first controls coefficients of pure powers using vanishing of first Koszul homology. A faithfully flat polynomial extension and a generic change of basis reduce all homogeneous coefficient relations to that special case, yielding the associated-graded conclusion..

\noindent Editorial provenance note (not author text). History: excerpt from revision \texttt{a0444\allowbreak 6e57e\allowbreak c1fbc\allowbreak 252a8\allowbreak 71afc\allowbreak ec775\allowbreak 2fb28\allowbreak 07b14}, more-algebra.tex, lines 7359--7438. Reference presentation was adapted; original-author context and core support blocks were retained or added. The mathematical wording is unchanged. Licensed under GNU FDL 1.2 or later; no invariant sections or cover texts. See ../licenses/COPYING. Context blocks reproduce author definitions and supporting results; they are not additional counted problems.

\begin{definition}
\label{definition-koszul-regular-sequence}
Let $R$ be a ring. Let $r \geq 0$ and let $f_1, \ldots, f_r \in R$
be a sequence of elements. Let $M$ be an $R$-module.
The sequence $f_1, \ldots, f_r$ is called
\begin{enumerate}
\item {\it $M$-Koszul-regular} if
$H_i(K_\bullet(f_1, \ldots, f_r) \otimes_R M) = 0$ for
all $i \not = 0$,
\item {\it $M$-$H_1$-regular} if
$H_1(K_\bullet(f_1, \ldots, f_r) \otimes_R M) = 0$,
\item {\it Koszul-regular} if $H_i(K_\bullet(f_1, \ldots, f_r)) = 0$ for
all $i \not = 0$, and
\item {\it $H_1$-regular} if $H_1(K_\bullet(f_1, \ldots, f_r)) = 0$.
\end{enumerate}
\end{definition}

\begin{lemma}
\label{lemma-H1-regular-quasi-regular}
An $M$-$H_1$-regular sequence is $M$-quasi-regular.
\end{lemma}

\begin{proof}
Let $R$ be a ring and let $M$ be an $R$-module.
Let $f_1, \ldots, f_r$ be an $M$-$H_1$-regular sequence.
Denote $J = (f_1, \ldots, f_r)$. The assumption means that we have
an exact sequence
$$
\wedge^2(R^r) \otimes M \to R^{\oplus r} \otimes M \to JM \to 0
$$
where the first arrow is given by
$e_i \wedge e_j \otimes m \mapsto (f_ie_j - f_je_i) \otimes m$.
Tensoring the sequence with $R/J$ we see that
$$
JM/J^2M = (R/J)^{\oplus r} \otimes_R M = (M/JM)^{\oplus r}
$$
is a finite free module. To finish the proof we have to prove
for every $n \geq 2$ the following: if
$$
\xi = \sum\nolimits_{|I| = n, I = (i_1, \ldots, i_r)}
m_I f_1^{i_1} \ldots f_r^{i_r} \in J^{n + 1}M
$$
then $m_I \in JM$ for all $I$. In the next paragraph, we prove
$m_I \in JM$ for $I = (0, \ldots, 0, n)$ and in the last paragraph
we deduce the general case from this special case.

\medskip\noindent
Let $I = (0, \ldots, 0, n)$. Let $\xi$ be as above. We can write
$\xi = m_1 f_1 + \ldots + m_{r - 1}f_{r - 1} + m_I f_r^n$.
As we have assumed $\xi \in J^{n + 1}M$, we can also write
$\xi = \sum_{1 \leq i \leq j \leq r - 1} m_{ij}f_if_j +
\sum_{1 \leq i \leq r - 1}m'_i f_if_r^n + m'' f_r^{n + 1}$.
Then we see that
$$
\begin{matrix}
(m_1 - m_{11}f_1 - m'_1f_r^n)f_1 + \\
(m_2 - m_{12}f_1 - m_{22}f_2 - m'_2f_r^n)f_2 + \\
\ldots + \\
(m_{r - 1} - m_{1 r - 1}f_1 - \ldots - m_{r - 1 r - 1}f_{r - 1}
- m'_{r - 1}f_r^n)f_{r - 1} + \\
(m_I - m'' f_r)f_r^n = 0
\end{matrix}
$$
Since $f_1, \ldots, f_{r - 1}, f_r^n$ is $M$-$H_1$-regular by
Lemma \href{https://stacks.math.columbia.edu/tag/062G}{lemma mult koszul regular (Tag 062G)} we see that $m_I - m'' f_r$
is in the submodule $f_1M + \ldots + f_{r - 1}M + f_r^nM$.
Thus $m_I \in f_1M + \ldots + f_rM$.

\medskip\noindent
Let $S = R[x_1, x_2, \ldots, x_r, 1/x_r]$. The ring map $R \to S$ is faithfully
flat, hence $f_1, \ldots, f_r$ is an $M$-$H_1$-regular sequence in $S$, see
Lemma \href{https://stacks.math.columbia.edu/tag/062H}{lemma koszul regular flat base change (Tag 062H)}.
By
Lemma \href{https://stacks.math.columbia.edu/tag/0625}{lemma change basis (Tag 0625)}
we see that
$$
g_1 = f_1 - \frac{x_1}{x_r} f_r,
\ \ldots,
\ g_{r - 1} = f_{r - 1} - \frac{x_{r - 1}}{x_r} f_r,
\ g_r = \frac{1}{x_r}f_r
$$
is an $M$-$H_1$-regular sequence in $S$. Finally, note that our element
$\xi$ can be rewritten
$$
\xi = \sum\nolimits_{|I| = n, I = (i_1, \ldots, i_r)}
m_I (g_1 + x_1 g_r)^{i_1} \ldots (g_{r - 1} + x_{r - 1} g_r)^{i_{r - 1}}
(x_rg_r)^{i_r}
$$
and the coefficient of $g_r^n$ in this expression is
$$
\sum m_I x_1^{i_1} \ldots x_r^{i_r}
$$
By the case discussed in the previous paragraph this sum is in
$J(M \otimes_R S)$. Since the monomials $x_1^{i_1} \ldots x_r^{i_r}$
form part of an $R$-basis of $S$ over $R$ we conclude that $m_I \in J$
for all $I$ as desired.
\end{proof}
\end{document}
